\clearpage
\SourceLeaf{356}
\pdfbookmark[2]{8. Gauss an Olbers, 29. Mai 1815}{parabolisch-008}
\begin{center}
[8.]\par
\medskip
\textsc{Gauss} an \textsc{Olbers}, Göttingen, 29. Mai 1815.
\end{center}
\noindent\hspace*{\parindent}.\ .\ .\ .\ .\ .\quad Meine theoretischen Untersuchungen über die Berechnung der Cometenbahnen im allgemeinen hätte ich wohl einige Neigung, in Zukunft einmal in einem eigenen Werke bekannt zu machen, als Supplement zu meiner Th. M. C. C. Es könnte vielleicht 6—8 Bogen stark werden, und ich würde dann noch eine Tafel für die parabolische Bewegung von einer neuen Einrichtung beifügen, deren Gebrauch noch etwas bequemer ist als der der \textsc{Barker}schen. Diese möchte auch noch 3 Bogen betragen. Der Grund dazu ist schon gelegt, wobei mir Hr. \textsc{Encke} noch geholfen hat.\quad .\ .\ .\ .\ .
\par\medskip
\begin{center}\rule{.18\textwidth}{.4pt}\end{center}
\medskip
\pdfbookmark[2]{9. Gauss an Bessel, 24. Junius 1815}{parabolisch-009}
\begin{center}
[9.]\par
\medskip
\textsc{Gauss} an \textsc{Bessel}, Göttingen, 24. Junius 1815.
\end{center}
\noindent\hspace*{\parindent}.\ .\ .\ .\ .\ .\quad Ich habe halb und halb die Absicht, demnächst ein Supplement der Theoria motus corporum coelestium für die parabolischen Bahnen zu geben. Ich habe dazu mehrere nicht ganz unerhebliche theoretische Zusätze; auch denke ich eine neue Tafel für die parabolische Bewegung beizufügen, noch bequemer eingerichtet als die \textsc{Barker}sche, deren Berechnung bereits so weit gediehen ist, dass ich mich schon jetzt zu meinem eigenen Gebrauch keiner andern mehr bediene.\quad .\ .\ .\ .\ .
\par\medskip
\begin{center}\rule{.22\textwidth}{.4pt}\end{center}
